\subsection{Biological Vulnerability and Intertemporal Health Investment}

Intertemporal preferences for vaccination timing are best understood not as abstract measures of impatience, but as responses shaped by biological vulnerability and exposure risk. When individuals face elevated baseline health risks, each day of delayed protection carries a greater expected biological cost. Steeper short-run discounting, in this light, reflects an adaptive valuation of immediacy under uncertainty---a rational response to risk rather than a departure from optimal health investment \parencite{Grossman1972, Attema2018}.

This interpretation finds support in the elevated delay sensitivity observed among respondents with chronic medical conditions. Chronic illness---whether cardiovascular, metabolic, or respiratory---compromises immune function, diminishes physiological reserves, and heightens susceptibility to severe infection outcomes. For these individuals, remaining unprotected is not merely inconvenient; it represents a period of amplified biological exposure. From a health-capital perspective, poorer baseline health steepens the risk gradient over time, making delayed vaccination disproportionately costly. The stronger preference for immediacy in this group thus reflects risk-weighted decision-making rather than generalized behavioural bias \parencite{Bok2024, Brewer2017}.

Age-related patterns reinforce this logic. Both younger and older adults exhibit steeper discounting than middle-aged respondents, but for distinct biological and behavioural reasons. Among older adults, immunosenescence---the progressive decline in immune competence with age---elevates both infection risk and disease severity, providing a clear physiological basis for urgency. For younger adults, steeper discounting may instead reflect greater uncertainty about future behavioural follow-through or competing life demands. In both cases, time preferences track perceived vulnerability across the life course rather than reflecting uniform impatience \parencite{Lawless2013, Wang2018}.

Institutional trust introduces a further layer of complexity. Lower trust heightens uncertainty about whether delayed vaccination will ultimately deliver its promised benefits---whether due to concerns about vaccine availability, quality, or institutional follow-through. Under such uncertainty, discounting future protection more steeply becomes a coherent response: if the probability of receiving effective protection diminishes, the rational weight placed on that future benefit declines accordingly. Trust, in this framework, functions as intertemporal capital---stabilising expectations about future health returns and thereby attenuating present bias \parencite{Brewer2017, Betsch2018}.

Taken together, these findings suggest that heterogeneity in vaccination timing preferences reflects adaptive responses to biological and institutional risk rather than idiosyncratic variation in patience. Viewing delay aversion through this lens clarifies why identical waiting periods impose unequal burdens across populations: for the chronically ill, the elderly, and those lacking institutional trust, each day of delay carries greater weight. This perspective underscores the importance of integrating biological vulnerability into analyses of intertemporal health behaviour and challenges interpretations that treat time preferences as stable, context-independent traits.
